\listfiles
\documentclass[\OPTclass]{book}
\usepackage[\OPTgeometry]{geometry}
\usepackage{amsmath,amsthm,amssymb}
\usepackage[utf8]{inputenc}
\usepackage[T1]{fontenc}
\usepackage[french]{babel}
\frenchbsetup{
  og=«, fg=»,
  ThinColonSpace=true
}
\addto\captionsfrench{%
  \renewcommand{\contentsname}{Sommaire}%
}
\usepackage[upright,widespace]{fourier}
\usepackage{maths}
\usepackage{microtype}
\usepackage[shortlabels]{enumitem}
\usepackage{booktabs}
\usepackage[x11names]{xcolor}
\usepackage{tikz}
\usepackage{graphicx}
\usepackage{titlesec}
\usepackage{titletoc}
\usepackage{fancyhdr}

\usepackage[
  colorlinks=true,
  linkcolor=Blue4,
  urlcolor=Blue4,
  linktocpage=true,
  pdftitle={Éléments de théorie des groupes. Solutions des exercices.},
  pdfauthor={Éric Guirbal},
  pdfsubject={Mathématiques},
  pdfkeywords={théorie des groupes, groupe monogène, groupe cyclique, groupe
    diédral, groupe fini, theorèmes de Sylow, sous-groupe normal, action de
    groupe, suite de composition, groupe abélien, groupe libre},
  pdflang=fr-FR]{hyperref}

%%% En-têtes
\pagestyle{fancy}
\fancyhf{}
\renewcommand{\chaptermark}[1]{%
  \markboth{\sffamily#1}{}}
\fancyhead[RE]{\sffamily Éléments de théorie des groupes. Solutions des exercices}
\fancyhead[LO]{\sffamily Chapitre \thechapter. \leftmark}
\fancyhead[RO,LE]{\thepage}
\renewcommand{\headrulewidth}{0pt}

%%% Chapitres
\titleformat{\chapter}[display]
  {\filcenter}
  {\textls[150]{\MakeUppercase{Chapitre\,\, \thechapter}}}
  {0.5pc}  
  {\LARGE\itshape}
  []
\titlespacing{\chapter}
  {0pt}
  {0ex}
  {10ex}

\renewcommand{\thechapter}{%                                                    
  \ifnum\value{chapter}=1
    premier%
  \else
    \Roman{chapter}%
  \fi
}

%%% Sommaire
\titlecontents{chapter}
  [0pc] % left
  {\addvspace{3pt}} % above
  {{\scshape Chapitre \thecontentslabel} / \itshape} % before with label
  {\scshape \thecontentslabel} % before without label
  {\titlerule*[0.5pc]{.}\contentspage} % filler and page
  [] % after

%%% Listes
\setlist[enumerate]{label=\alph*),leftmargin=*,widest=a}

%%% Environnements
\theoremstyle{definition}
\newtheorem{solution}{}
\theoremstyle{remark}
\newtheorem*{remarque}{Remarque}

%%% Pour que les fausses pages vides (de fin de chapitre par exemple)
%%% soient complètement vides (sans en-tête, pied de page ou folio)
\makeatletter
\renewcommand*\cleardoublepage{
  \clearpage
  \if@twoside
    \ifodd
      \c@page
    \else
      \hbox{}
      \thispagestyle{empty}
      \newpage
      \if@twocolumn
        \hbox{}
        \newpage
      \fi
    \fi
  \fi
}
\makeatother

%%% Le dossier contenant les solutions
\newcommand*{\solutionsdir}{solutions}

%%% Pour chacun des chapitres, le nombre d'exercices
%%% et le dossier dans lequel sont leurs solutions.
\newcommand*{\thebook}{%
  % dossier/chapitre/nombre d'exercices
  1-structure_de_groupe/Structure de groupe/33,
  2-classes_modulo_un_sous_groupe/Classes modulo un sous-groupe/10,
  3-groupes_monogenes/{Groupes monogènes, symétriques et diédraux}/33,
  4-sous_groupes_normaux/Sous-groupes normaux/38,
  5-groupe_operant_sur_un_ensemble/Groupe opérant sur un ensemble/27,
  6-groupes_finis/Groupes finis. Théorèmes de Sylow/26,
  7-suites_de_composition/Suites de composition/35,
  8-groupes_abeliens/Groupes abéliens/23,
  9-groupes_libres/Groupes libres. Générateurs et relations. Produit libre de groupes/12%
}
%%% Détermine les nombres des solutions pour chacun des chapitres
%%% ansi que les nombres totaux des exercices et solutions.
%%% Ces nombres sont sauvegardés dans des compteurs.
\newcommand*{\calculatewhathasbeendone}{
  \newcounter{totalexercises}
  \newcounter{totalsolutions}
  \foreach \dir/\chaptertitle/\numberofexercises [count=\chapternum] in \thebook {%
    \newcounter{solutionsforchapter\chapternum}%
    \foreach \sol in {1,...,\numberofexercises} {%
      \edef\solutionfilename{%
        \ifnum \sol < 10%
          \solutionsdir/\dir/\chapternum-0\sol.tex%
        \else%
          \solutionsdir/\dir/\chapternum-\sol.tex%
        \fi%
      }%
      \IfFileExists{\solutionfilename}{%
        \stepcounter{solutionsforchapter\chapternum}%
      }{}%
    }%
    \addtocounter{totalexercises}{\numberofexercises}%
    \addtocounter{totalsolutions}{\value{solutionsforchapter\chapternum}}%
  }
}
%%% Représente un pourcentage par un donut
\newcommand*{\donut}[8]{%
  \begin{tikzpicture}
    \pgfmathtruncatemacro{\smallquantity}{#1}
    \pgfmathtruncatemacro{\largequantity}{#2}
    \pgfmathsetmacro{\radint}{#3} % rayon intérieur
    \pgfmathsetmacro{\radext}{#4} % rayon extérieur
    \def\colorint{#5}             % couleur du disque intérieur
    \def\colorext{#6}             % couleur du donut
    \def\colorbrd{#7}             % couleur des bords du donut
    \def\colordonut{#8}           % couleur de la part du donut qui représente p.100
    \pgfmathtruncatemacro{\percent}{round(100*\smallquantity/\largequantity}

    \fill[\colorint] circle(\radint);
    \filldraw[even odd rule,fill=\colorext,draw=\colorbrd] 
      circle(\radext) circle(\radint) 
      node[\colordonut,scale=4] {$\percent\,\%$} 
      node[\colordonut,yshift=-0.85cm,scale=1.2] {$\smallquantity/\largequantity$};
    
    \pgfmathsetmacro{\angle}{3.6*\percent};
    \fill[\colordonut] (0,\radint) arc (90:90-\angle:\radint) 
      -- (90-\angle:\radext) arc (90-\angle:90:\radext) 
      -- cycle;
  \end{tikzpicture}
}
%%% Représente un pourcentage
\newcommand*{\pie}[6]{%
  \begin{tikzpicture}[baseline=-#3cm/2.5]
    \pgfmathsetmacro{\smallquantity}{#1}
    \pgfmathsetmacro{\largequantity}{#2}
    \pgfmathsetmacro{\radius}{#3}
    \def\colorpie{#4}   % couleur de la part représentant le p.100
    \def\colorbg{#5}    % couleur de fond
    \def\colorbd{#6}    % couleur du bord
    \pgfmathsetmacro{\angle}{360*\smallquantity/\largequantity};
    
    \filldraw[fill=\colorbg,draw=\colorbd] circle(\radius);
    \fill[\colorpie] (0,\radius) arc (90:90-\angle:\radius) -- (0,0) -- cycle;
  \end{tikzpicture}%
}

\begin{document}
\title{
  Éléments de théorie des groupes\\
  Solutions des exercices
}

\author{Éric Guirbal}
\date{}

\frontmatter

\maketitle

%%%%%%%%%%%%%%%%%%% COPYRIGHT %%%%%%%%%%%%%%%%%%%%%%%%%%%%%%%%%%%%%%%%%%%%%%%%
\input{revision}
\thispagestyle{empty}
\null
\vfill
\begin{center}
  \small

  \textbf{%
    Éléments de théorie des groupes \\
    Solutions des exercices
  }
  \medskip

  Éric \textsc{Guirbal}

  \medskip

  Version: \gitDescribe{} (\gitDate)

  Compilé le \today

  \vfill

  \includegraphics[width=2cm]{images/cc-by-nc-sa} \\ [0.15cm]

  Ce document est distribué selon les termes de la licence Creative Commons
  Attribution - Pas d'utilisation commerciale - Partage à l'identique 3.0
  France. \\
  \url{https://creativecommons.org/licenses/by-nc-sa/3.0/fr/}
\end{center}


\tableofcontents

\chapter{État d'avancement du projet}

\calculatewhathasbeendone

\vfill

%%% Pourcentage des exercices qui ont une solution
\begin{center}
  \donut
    {\value{totalsolutions}}
    {\value{totalexercises}}
    {2}         % rayon intérieur (en cm)
    {2.5}       % rayon extérieur (en cm)
    {white}     % couleur du disque intérieur
    {black!15}  % couleur du donut
    {black!5}   % couleur des bords du donut
    {black}     % couleur de la part du donut qui représente le p.100
\end{center}

\bigskip

%%% Pourcentage des exercices qui ont une solution pour chacun des chapitres
{
  \footnotesize
  \begin{enumerate}[I.]
    \foreach \dir / \chaptertitle / \numberofexercises [count=\chapternum] in \thebook {%
      \item
        \begin{minipage}[t]{0.84\linewidth}
          \chaptertitle~~\titlerule*[0.5pc]{.}
        \end{minipage}%
        \hfill\arabic{solutionsforchapter\chapternum}/\numberofexercises\hspace{0.2cm}
        \pie
          {\value{solutionsforchapter\chapternum}}
          {\numberofexercises}
          {0.2}     % rayon (en cm)
          {black}   % couleur de la part
          {white}   % couleur de fond
          {black}   % couleur du bord
       }%
  \end{enumerate}
}

\vfill\vfill


\mainmatter

\foreach \dir / \chaptertitle / \numberofexercises [count=\chapternum] in \thebook {
  \chapter{\chaptertitle}
  \foreach \sol in {1,...,\numberofexercises} {%
    \edef\solutionfilename{
      \ifnum \sol < 10
        \solutionsdir/\dir/\chapternum-0\sol.tex
      \else
        \solutionsdir/\dir/\chapternum-\sol.tex
      \fi
    }
    \IfFileExists{\solutionfilename}{%
      \setcounter{solution}{\numexpr\sol-1\relax}
      \begin{solution}
        \input{\solutionfilename}
      \end{solution}
    }{}
  }
}

\backmatter

\bibliographystyle{alpha}
\bibliography{bibliographie}
\nocite{*}
\end{document}

